% Complete finite-filling certificates, retaining the original marked periods.
% Host prerequisites: candidate_geometry.tex, exact_algebra.tex;
% normal-line comparison: analytic_picard.tex. No package or macro definitions.
\subsection{Exact product covers, normalization branches, and peripheral kernels of the finite fillings}
\label{ffcert:sec:finite-fillings}

We retain the actual period functions and finite fillings of
\cref{prop:period-input-audit,prop:monodromy-certificate,%
prop:candidate-finite-fillings}. The following explicit product maps
describe the original central tori and surfaces, not replacements of
their periods. We also identify the complete varying family as the
normalization of the actual ramified base change, compute the entire
peripheral inclusion kernel, and give a smooth trivialization of the
finite normal line. These calculations use no sphere identification,
global integral cohomology, or CDP theorem. The separate period and cusp
constructions retain their own proofs and are not inferred from these
finite-local calculations.

\subsubsection{The original lattice, periods, and their complex summands}

Use the ordered marked lattice and its dual basis
\[
 \Lambda=\mathbb Z\langle\widehat\gamma,\widehat u,
            \widehat w,\widehat\delta\rangle,
 \qquad \Lambda^*=\mathbb Z\langle\gamma,u,w,\delta\rangle.
 \tag{FF1}\label{ffcert:eq:marking}
\]
In these original column coordinates the matrices and translations are
\[
\begin{gathered}
 A_1=\begin{pmatrix}1&0&0&0\\6&0&1&0\\-6&-1&-1&0\\-2&1&0&1\end{pmatrix},
 \qquad
 A_2=\begin{pmatrix}1&0&0&0\\0&0&-1&0\\-6&1&0&0\\3&0&1&1\end{pmatrix},\\
 \varepsilon=(1,2,-4,0)^{\mathsf t},\quad
 \varepsilon'=(1,3,-3,0)^{\mathsf t},\quad
 v_1=\varepsilon,\quad v_2=-\varepsilon',\quad (m_1,m_2)=(3,4).
\end{gathered}
\tag{FF2}\label{ffcert:eq:matrices-translations}
\]
They satisfy $A_j^{m_j}=I$, $A_jv_j=v_j$, and
$\gamma A_j=\gamma$, with $\gamma(v_1)=1$ and $\gamma(v_2)=-1$.
Retain the entire period matrix
\[
 \Pi(z)=\begin{pmatrix}6\mu(z)&\tau(z)&1&0\\
                         \beta(z)&\mu(z)&0&1\end{pmatrix},
 \qquad
 D(z)=\operatorname{Im}\beta(z)
       -\frac{6(\operatorname{Im}\mu(z))^2}{\operatorname{Im}\tau(z)}<0.
 \tag{FF3}\label{ffcert:eq:periods}
\]
At the original elliptic centres let
\[
\begin{gathered}
 \rho=e^{\pi i/3},\quad \tau(z_1)=\rho,\quad \tau(z_2)=i,
 \quad \mu_1:=\mu(z_1)=\frac{2-\rho}{3},\quad
 \mu_2:=\mu(z_2)=\frac{1-i}{2},\\
 \beta_j=\beta(z_j),\qquad
 \kappa_1=\beta_1+2\mu_1,\qquad
 \kappa_2=\beta_2+3\mu_2.
\end{gathered}
\tag{FF4}\label{ffcert:eq:central-parameters}
\]
Neither $\beta_j$ nor either of its real components is adjusted.
Since $\operatorname{Im}\mu_1=-\operatorname{Im}\rho/3$ and
$\operatorname{Im}\mu_2=-1/2$, the exact inequalities are
\[
 \operatorname{Im}\kappa_1
 =\operatorname{Im}\beta_1-\tfrac23\operatorname{Im}\rho=D(z_1)<0,
 \qquad
 \operatorname{Im}\kappa_2
 =\operatorname{Im}\beta_2-\tfrac32=D(z_2)<0.
 \tag{FF5}\label{ffcert:eq:negative-orientation}
\]
Thus $(1,\kappa_j)$ is real-linearly independent with its original
negative imaginary orientation. No conjugation or positive-imaginary
replacement is made.

Put $\delta_*:=\widehat\delta$ and define
\[
\begin{aligned}
 e_1&=(0,1,-2,0)^{\mathsf t},&f_1&=(0,0,-3,1)^{\mathsf t},\\
 e_2&=(0,1,-1,0)^{\mathsf t},&f_2&=(0,0,-2,1)^{\mathsf t}.
\end{aligned}
\tag{FF6}\label{ffcert:eq:moving-generators}
\]

\begin{proposition}[The full fixed and moving integral lattices]
\label{ffcert:prop:integral-summands}
The exact fixed and image lattices are
\[
\begin{aligned}
 \Lambda^{A_1}&=\mathbb Z\varepsilon\oplus\mathbb Z\delta_*,&
 (A_1-I)\Lambda&=\mathbb Ze_1\oplus\mathbb Zf_1,\\
 \Lambda^{A_2}&=\mathbb Z\varepsilon'\oplus\mathbb Z\delta_*,&
 (A_2-I)\Lambda&=\mathbb Ze_2\oplus\mathbb Zf_2.
\end{aligned}
\tag{FF7}\label{ffcert:eq:integral-summands}
\]
The image lattices are saturated, with complete annihilators
\[
\begin{gathered}
 \operatorname{Ann}((A_1-I)\Lambda)
 =\mathbb Z\gamma\oplus\mathbb Z\psi_1,
 \quad \psi_1=2u+w+3\delta,\\
 \operatorname{Ann}((A_2-I)\Lambda)
 =\mathbb Z\gamma\oplus\mathbb Z\psi_2,
 \quad \psi_2=u+w+2\delta,\\
 \det[\varepsilon,\delta_*,e_1,f_1]=-3,
 \qquad \det[\varepsilon',\delta_*,e_2,f_2]=-2.
\end{gathered}
\tag{FF8}\label{ffcert:eq:annihilators-indices}
\]
In particular the sum of the fixed and moving lattices is not being
asserted to be all of $\Lambda$.
\end{proposition}

\begin{proof}
Solving $(A_1-I)(a,b,c,d)^{\mathsf t}=0$ gives $b=2a,c=-4a$;
the second equation gives $b=3a,c=-3a$. In both cases $a,d$ are arbitrary
integers. The columns of the two image matrices are, in order,
\[
 \begin{aligned}
 A_1-I &: \quad 6e_1-2f_1,\ -e_1+f_1,\ e_1,\ 0,\\
 A_2-I &: \quad 3f_2,\ -e_2,\ -e_2+f_2,\ 0.
 \end{aligned}
 \tag{FF9}\label{ffcert:eq:image-columns}
\]
These columns generate precisely the two stated image lattices. The first
is exactly the set of integer columns satisfying $a=0$ and
$2b+c+3d=0$: such a column is $be_1+df_1$. The second is exactly
the set satisfying $a=0$ and $b+c+2d=0$, with the analogous expression
$be_2+df_2$. They are therefore saturated. Their annihilators are the
displayed full integer spans: in the first case an annihilating row has
coefficients $(a,2c,c,3c)$, and in the second $(a,c,c,2c)$.
Expansion of the two four-by-four determinants gives $-3$ and $-2$.
\end{proof}

The complex images under the unchanged periods are
\[
\begin{aligned}
 \Pi(z_1)\varepsilon&=(0,\kappa_1),&
 \Pi(z_2)\varepsilon'&=(0,\kappa_2),&
 \Pi(z_j)\delta_*&=(0,1),\\
 \Pi(z_1)e_1&=(\rho-2,-(\rho-2)/3),&
 \Pi(z_1)f_1&=(-3,1),\\
 \Pi(z_2)e_2&=(i-1,-(i-1)/2),&
 \Pi(z_2)f_2&=(-2,1).
\end{aligned}
\tag{FF10}\label{ffcert:eq:complex-images}
\]
For example the first fixed coordinate vanishes because
$6(2-\rho)/3+2\rho-4=0$ or $6(1-i)/2+3i-3=0$, respectively.
The fixed real plane consequently maps onto the complex line
$\{(0,\xi):\xi\in\mathbb C\}$. The moving real planes map onto
$\{(\eta,-\eta/3)\}$ and $\{(\eta,-\eta/2)\}$, respectively,
because their displayed period pairs are real-linearly independent.
These are complementary complex lines in the original $\mathbb C^2$.

\subsubsection{The exact elliptic-product covers of the central tori and surfaces}

Define the one-dimensional complex tori
\[
\begin{gathered}
 E_j=\mathbb C/(\mathbb Z+\kappa_j\mathbb Z),\\
 F_1=\mathbb C/(\mathbb Z(\rho-2)+\mathbb Z(-3)),\qquad
 F_2=\mathbb C/(\mathbb Z(i-1)+\mathbb Z(-2)).
\end{gathered}
\tag{FF11}\label{ffcert:eq:elliptic-factors}
\]
Each period pair is real-linearly independent. Translation by its rank-two
lattice is free and properly discontinuous; the image of a closed real
parallelotope covers the quotient, proving compactness. Thus these are
smooth elliptic curves with the class of zero distinguished.

\begin{theorem}[Product maps onto the original finite central surfaces]
\label{ffcert:thm:product-covers}
The complex-linear maps
\[
 L_1(\xi,\eta)=(\eta,\xi-\eta/3),\qquad
 L_2(\xi,\eta)=(\eta,\xi-\eta/2)
 \tag{FF12}\label{ffcert:eq:addition-maps}
\]
induce surjective local biholomorphic group homomorphisms
\[
 \ell_j:E_j\times F_j\longrightarrow
 A_0^{(j)}:=\mathbb C^2/\Pi(z_j)\Lambda,
 \quad \ker\ell_1=\langle k_1\rangle,
 \quad \ker\ell_2=\langle k_2\rangle,
 \tag{FF13}\label{ffcert:eq:torus-covers}
\]
where
\[
 k_1=(2/3,\rho),\quad \operatorname{ord}(k_1)=3,
 \qquad k_2=(1/2,i),\quad \operatorname{ord}(k_2)=2.
 \tag{FF14}\label{ffcert:eq:kernel-points}
\]
The following automorphisms lift the original affine generators:
\[
 g^{\mathrm{prod}}_1(\xi,\eta)
   =(\xi+\kappa_1/3,(\rho-1)\eta),\qquad
 g^{\mathrm{prod}}_2(\xi,\eta)
   =(\xi-\kappa_2/4,-i\eta).
 \tag{FF15}\label{ffcert:eq:product-generators}
\]
Writing $k_j$ also for translation by that point, their joint groups
are $\langle k_1,g^{\mathrm{prod}}_1\rangle\cong C_3\times C_3$
and $\langle k_2,g^{\mathrm{prod}}_2\rangle\cong C_2\times C_4$.
They act freely, and the original reduced central surfaces are exactly
\[
 S_1=(E_1\times F_1)/\langle k_1,g^{\mathrm{prod}}_1\rangle,
 \qquad
 S_2=(E_2\times F_2)/\langle k_2,g^{\mathrm{prod}}_2\rangle.
 \tag{FF16}\label{ffcert:eq:surface-product-covers}
\]
The finite \'etale product-cover degrees are $9$ and $8$.
\end{theorem}

\begin{proof}
Each $L_j$ has determinant $-1$ in the displayed product and original
$(\zeta_1,\zeta_2)$ coordinates. Equation~\eqref{ffcert:eq:complex-images}
shows that the product period lattice maps exactly to the sum of the
fixed and moving sublattices inside $\Pi(z_j)\Lambda$.
Thus $L_j$ induces the stated surjective local biholomorphism, and its
kernel is identified with
$\Lambda/(\Lambda^{A_j}+(A_j-I)\Lambda)$, of order $3$ or $2$ by
\eqref{ffcert:eq:annihilators-indices}. Explicitly,
\[
 \widehat u=\tfrac23\delta_*+e_1-\tfrac23f_1,
 \qquad
 \widehat u=\tfrac12\delta_*+e_2-\tfrac12f_2.
 \tag{FF17}\label{ffcert:eq:kernel-lifts}
\]
Consequently $L_j(k_j)=\Pi(z_j)\widehat u$. The first kernel point
has order exactly three, since its fixed coordinate does, and
$3\rho=3(\rho-2)-2(-3)$ is a moving period. The second has order
exactly two, since its fixed coordinate does, and
$2i=2(i-1)-(-2)$ is a moving period. These orders equal the full kernel
cardinalities, proving the complete kernel assertions.

At the centres the original covariance matrices are
\[
 R_1=\begin{pmatrix}\rho-1&0\\(2-\rho)/3&1\end{pmatrix},\qquad
 R_2=\begin{pmatrix}-i&0\\(1+i)/2&1\end{pmatrix}.
 \tag{FF18}\label{ffcert:eq:central-covariance}
\]
They act by $1$ on the fixed lines and by
$\rho-1=e^{2\pi i/3}$ or $-i$ on the corresponding moving line.
Substitution gives the full equality of affine maps on $\mathbb C^2$,
\[
 L_jg_j^{\mathrm{prod}}
   =R_jL_j+\Pi(z_j)v_j/m_j.
 \tag{FF19}\label{ffcert:eq:product-intertwining}
\]
In particular the negative translation $-\kappa_2/4$ is forced by
$v_2=-\varepsilon'$. It has not been replaced by its opposite.

The moving rotations preserve their exact period lattices. For example,
using $\rho^2=\rho-1$, their images of the two period generators are
\[
\begin{aligned}
 (\rho-1)(\rho-2)&=-2(\rho-2)+(-3),&
 (\rho-1)(-3)&=-3(\rho-2)+(-3),\\
 (-i)(i-1)&=(i-1)-(-2),&
 (-i)(-2)&=2(i-1)-(-2).
\end{aligned}
\tag{FF20}\label{ffcert:eq:moving-rotations}
\]
The first matrix in these period bases has columns $(-2,1)$ and
$(-3,1)$, determinant $1$; the second has columns $(1,-1)$ and
$(2,-1)$, determinant $1$.
Equivalently, invariance and invertibility follow by applying the finite
orders of the rotations to these integral period images.
The kernel classes are fixed by the moving rotations because
\[
 (\rho-1)\rho-\rho=-1-\rho=-(\rho-2)+(-3),\qquad
 (-i)i-i=1-i=-(i-1).
 \tag{FF21}\label{ffcert:eq:kernel-rotation}
\]
Hence the product automorphism commutes with translation by $k_j$ on
the product torus. It has exact order $m_j$: its moving rotation has
that order; its $m_j$-th fixed-coordinate translation is
$\kappa_1$ or $-\kappa_2$; and no smaller positive multiple of that
translation is a fixed-coordinate period, since $\kappa_j$ is nonreal.

To prove both the direct-product assertion and freeness of the entire
joint group, put $(d_1,d_2)=(3,2)$, $(c_1,c_2)=(2/3,1/2)$ and
$(\epsilon_1,\epsilon_2)=(1,-1)$. A fixed point for $k_j^a
(g_j^{\mathrm{prod}})^b$, with $0\le a<d_j$ and $0\le b<m_j$,
would require
\[
 ac_j+\epsilon_jb\kappa_j/m_j=n+q\kappa_j,
 \qquad n,q\in\mathbb Z.
 \tag{FF22}\label{ffcert:eq:joint-freeness}
\]
Imaginary parts give $\epsilon_jb/m_j=q$, hence $b=0$. The real
equality then gives $ac_j\in\mathbb Z$, hence $a=0$ in its stated
range. Every nonidentity pair is therefore fixed-point free; in
particular it is not the identity transformation. This proves the
claimed joint groups. Quotienting first by $k_j$ and then by the lifted
affine generator gives exactly the original $S_j$ of
\cref{prop:candidate-finite-fillings}, by
\eqref{ffcert:eq:product-intertwining}. The total covering degrees are
$d_jm_j=9,8$ and the actions are free, so the covers are \'etale.
\end{proof}

\subsubsection{The varying quotient and its full local analytic coordinates}

Retain the original linearizing coordinate $s_j$ on $\Delta_j$, with
its actual round-disc radius $r_j>0$, and the original phase
\[
 s_j(g_jz)=\zeta_js_j(z),\quad
 \zeta_j=e^{-2\pi i/m_j},\quad t_j=s_j^{m_j},
 \qquad |s_j|<r_j,\quad |t_j|<r_j^{m_j}.
 \tag{FF23}\label{ffcert:eq:disc-and-phase}
\]
The radius is not replaced by one. On the universal marked cover the
holomorphic affine map and the exact period covariance are
\[
\begin{gathered}
 \widehat G_j(z,\zeta)
  =\bigl(g_jz,R_{g_j}(z)\zeta+\Pi(g_jz)v_j/m_j\bigr),\\
 R_{g_j}(z)\Pi(z)=\Pi(g_jz)A_j,
 \qquad t_\lambda(z,\zeta)=(z,\zeta+\Pi(z)\lambda).
\end{gathered}
\tag{FF24}\label{ffcert:eq:varying-action}
\]
Let $G_j$ denote the induced map of the original varying torus family
$J'_j=(\Delta_j\times\mathbb C^2)/\Lambda$ and put
$N_j=J'_j/\langle G_j\rangle$, with quotient map $q_j$.
The real-analytic marking $\Xi(z,x)=(z,\Pi(z)x)$ converts its powers
to the exact formulas
\[
 \Xi^{-1}\widehat G_j^k\Xi(z,x)
   =(g_j^kz,A_j^kx+kv_j/m_j),\qquad
 \widehat G_j^{m_j}=t_{v_j}.
 \tag{FF25}\label{ffcert:eq:varying-powers}
\]

\begin{proposition}[Smooth quotient charts and the exact multiplicities]
\label{ffcert:prop:quotient-charts}
The map $q_j:J'_j\to N_j$ is a finite \'etale covering of degree
$m_j$. The quotient is a Hausdorff smooth complex threefold, and its
proper map $f_j:N_j\to D_j$ has exact scheme fibre $m_jS_j$ at the
centre. In its local quotient coordinates the transverse map is
literally $t_j=s_j^{m_j}$.
\end{proposition}

\begin{proof}
At $s_j=0$ the freeness is proved in
\cref{ffcert:thm:product-covers}. Away from zero, a power with
$0<k<m_j$ has base coordinate multiplied by $\zeta_j^k\ne1$, so
has no fixed point. Equation~\eqref{ffcert:eq:varying-powers} shows
that the $m_j$-th power is the original lattice translation, hence is
the identity on $J'_j$. This gives a free action of exact order $m_j$.

For $y\in J'_j$, its finitely many distinct orbit points can be
separated by pairwise disjoint coordinate neighbourhoods. Intersecting
their inverse translates and shrinking gives a coordinate neighbourhood
$U$ of $y$ with the $G_j^aU$ pairwise disjoint. Then
$q_j|_U:U\to q_j(U)$ is a biholomorphism. Two distinct finite orbits
have disjoint saturated neighbourhoods by the same finite separation
argument. These charts prove the Hausdorff and smooth assertions, and
the quotient is a finite \'etale covering. They give directly the
trivial-stabilizer case of the quotient analytic structure in
\cite[conditions (I)--(II), Theorem~1]{Cartan1954Quotient}; smoothness
is not inferred from a theorem about arbitrary finite quotients.

At a point of $S_j$, use the quotient chart to retain $s_j$ as a
transverse coordinate, together with two fibre coordinates. The full
map is $t_j=s_j^{m_j}$. Thus $S_j$ is a smooth reduced Cartier
divisor, and its scheme multiplicity is exactly $m_j$.

For compact $C\subset\Delta_j$, the map induced by $\Xi$ from
$C\times[0,1]^4$ onto the inverse image of $C$ in $J'_j$ is
continuous and surjective, so that inverse image is compact. The actual
power map between the discs of radii $r_j$ and $r_j^{m_j}$ is proper:
the inverse image of a compact target set is closed, bounded, and stays
strictly inside the radius-$r_j$ disc. For compact $K\subset D_j$,
the inverse image in $J'_j$ of all roots of its $t_j$-values is
therefore compact. Its finite quotient image is exactly $f_j^{-1}(K)$,
proving properness of the full restriction.
\end{proof}

\subsubsection{Every branch in the normalization of the actual base change}

Let $\Delta_j^{\mathrm{base}}$ be the same radius-$r_j$ disc with
coordinate $u$, mapped to $D_j$ by $t_j=u^{m_j}$, and form
\[
 P_j=N_j\times_{D_j}\Delta_j^{\mathrm{base}},\qquad
 \nu_j:J'_j\longrightarrow P_j,\quad
 \nu_j(y)=(q_j(y),s_j(y)).
 \tag{FF26}\label{ffcert:eq:normalization-map}
\]

\begin{theorem}[The original varying family is the normalization]
\label{ffcert:thm:normalization}
The map $\nu_j$ is exactly the normalization of $P_j$. It is a
biholomorphism away from $u=0$; at the centre it has all $m_j$ local
branches prescribed by the original cyclic action. Its central fibre
is the connected original torus $A_0^{(j)}$, and
\[
 \nu_j(G_jy)=(q_j(y),\zeta_js_j(y)).
 \tag{FF27}\label{ffcert:eq:normalization-deck}
\]
Thus the original negative-phase action lifts $u\mapsto\zeta_ju$,
permutes the local branches cyclically, and has quotient $N_j$.
\end{theorem}

\begin{proof}
The definition uses no logarithm choice. If $u\ne0$, the $s_j$-values
in a covering orbit are exactly the $m_j$ distinct roots of its
$t_j$-value, so there is exactly one representative with value $u$.
If $u=0$, all $m_j$ covering points have $s_j=0$. This proves
surjectivity and describes all fibres. A compact subset of $P_j$
projects to a compact subset of $N_j$, whose inverse image under
$q_j$ is compact. Its inverse image under $\nu_j$ is a closed subset
of that compact set, hence is compact. Thus $\nu_j$ is proper with
finite fibres. Away from the centre the unique orbit representative
and quotient charts give the holomorphic inverse.

More precisely, around $n\in S_j$ choose a chart $q_j(U)$ as in
\cref{ffcert:prop:quotient-charts}, with inherited transverse coordinate
$s$ and fibre coordinates $z_1,z_2$. The fibre product is locally
\[
 s^{m_j}-u^{m_j}
   =\prod_{a=0}^{m_j-1}(s-e^{2\pi ia/m_j}u)=0.
 \tag{FF28}\label{ffcert:eq:all-branches}
\]
The factors are distinct smooth prime hypersurface factors in the
regular convergent-power-series ring $\mathbb C\{z_1,z_2,s,u\}$.
Their product ideal is radical: in this factorial ring, divisibility
of a power by every distinct prime factor implies divisibility of the
element by their product. Thus the fibre-product germ is reduced.
Its total ring of fractions is the product of the branch fraction
fields. Its integral closure is the product of the regular branch
rings $\mathbb C\{z_1,z_2,u\}$: adjoining the branch idempotents is
integral, since each satisfies $e^2-e=0$, and each component is
integrally closed. This proves the complete local normalization,
including all central branches.

On $G_j^aU$ the second component of $\nu_j$ is $u=\zeta_j^a s$,
or equivalently
\[
 s=\zeta_j^{-a}u=e^{2\pi ia/m_j}u.
 \tag{FF29}\label{ffcert:eq:branch-identification}
\]
Each of the $m_j$ disjoint covering charts therefore maps
isomorphically onto exactly one branch of
\eqref{ffcert:eq:all-branches}; these exhaust the branches. The
local identifications glue by the specified map $\nu_j$, giving the
global normalization. The source is the original varying family,
not an isotrivial product. Restricting its coordinate to zero gives
its connected torus $A_0^{(j)}$. Equation~\eqref{ffcert:eq:normalization-deck}
follows directly from $s_j(G_jy)=\zeta_js_j(y)$ and $q_jG_j=q_j$.
\end{proof}

\subsubsection{Universal covers and the full peripheral inclusion kernel}

\begin{theorem}[Exact finite-filling deck groups and attachment kernel]
\label{ffcert:thm:peripheral-kernel}
The universal cover $\Delta_j\times\mathbb C^2\to N_j$ has the
following complete deck-group presentation:
\[
 \Gamma_j=\left\langle\Lambda,\widehat G_j\ \middle|\
 \begin{array}{l}
 \Lambda\text{ abelian},\\
 \widehat G_jt_\lambda\widehat G_j^{-1}=t_{A_j\lambda},\\
 \widehat G_j^{m_j}=t_{v_j}
 \end{array}\right\rangle.
 \tag{FF30}\label{ffcert:eq:full-deck-group}
\]
This is also the universal-cover group of $S_j$. On the punctured
filling the complete group is $\Lambda\rtimes_{A_j}\mathbb Z$,
with the generator specified below, and the entire inclusion kernel is
\[
 \ker\{\pi_1(N_j\setminus S_j)\to\pi_1(N_j)\}
   =\langle\widetilde G_j^{m_j}t_{-v_j}\rangle
   \cong\mathbb Z.
 \tag{FF31}\label{ffcert:eq:whole-kernel}
\]
It is a central subgroup, not merely a necessary power relation.
\end{theorem}

\begin{proof}
In the real marking the original universal-cover transformations are
\[
 t_\lambda(s,x)=(s,x+\lambda),\qquad
 \widehat G_j(s,x)=(\zeta_js,A_jx+v_j/m_j).
 \tag{FF32}\label{ffcert:eq:cover-real-action}
\]
They satisfy the two relations in
\eqref{ffcert:eq:full-deck-group} by $A_jv_j=v_j$.
Moving translations to the left and using the power relation expresses
every word as $t_\lambda\widehat G_j^k$ with $\lambda\in\Lambda$
and $0\le k<m_j$. These forms are unique: the base action determines
$k$, and the remaining translation determines $\lambda$. Therefore
the abstract presented group acts faithfully and no additional relation
is implicit.

For $0<k<m_j$, a fixed-point equation would imply, on applying the
original integral covector $\gamma$ to the real fibre coordinate,
\[
 \gamma(\lambda)+k\gamma(v_j)/m_j=0.
 \tag{FF33}\label{ffcert:eq:cover-freeness}
\]
This is impossible because $\gamma(\lambda)$ is integral and
$\gamma(v_j)=1$ or $-1$. A nonzero translation with $k=0$ has no
fixed point either. The lattice action is properly discontinuous:
on compact base sets the real inverses $\Pi(z)^{-1}$ are uniformly
bounded, leaving finitely many lattice translations that can meet two
fixed compact sets. The full group is a finite union of lattice cosets,
so the same property holds for it. Thus the simply connected
$\Delta_j\times\mathbb C^2$ is the universal cover. Its restriction
to $s=0$ is the simply connected $\mathbb C^2$ covering $S_j$ with
the same action and group.

On the punctured disc retain exactly
\[
 w=\frac{\log s_j}{2\pi i},\qquad s_j=e^{2\pi iw},\qquad
 \operatorname{Im}w> -\frac{\log r_j}{2\pi},
 \qquad
 \widetilde G_j(w,x)=(w-1/m_j,A_jx+v_j/m_j).
 \tag{FF34}\label{ffcert:eq:log-cover-action}
\]
The displayed half-plane times $\mathbb C^2$, with its pulled-back
holomorphic coordinates, is simply connected. Along with the lattice
translations, $\widetilde G_j$ generates its entire punctured deck
group: the ordinary logarithm deck map is already
\[
 Q(w,x)=(w+1,x)=t_{v_j}\widetilde G_j^{-m_j}(w,x).
 \tag{FF35}\label{ffcert:eq:ordinary-log-deck}
\]
There is no finite-power relation for $\widetilde G_j$, since its
$k$-th power has $w$-displacement $-k/m_j$. Thus the exact normal
forms are $t_\lambda\widetilde G_j^k$ for unrestricted
$k\in\mathbb Z$, and the group is precisely the stated semidirect
product with action $A_j$.

The inclusion sends $t_\lambda$ to $t_\lambda$ and
$\widetilde G_j$ to $\widehat G_j$. Write $k=qm_j+r$ with
$0\le r<m_j$. Its image has normal form
$t_{\lambda+qv_j}\widehat G_j^r$, which is the identity exactly
when $r=0$ and $\lambda=-qv_j$. This proves
\eqref{ffcert:eq:whole-kernel} in both directions. Its generator has
$w$-displacement $-1$, so has infinite order. It commutes with all
translations because $A_j^{m_j}=I$, and with
$\widetilde G_j$ because $A_jv_j=v_j$, proving centrality.
\end{proof}

\subsubsection{The original punctured overlap and integral cohomology pullback}

On the affine punctured model choose a local logarithm and put
\[
 \alpha(z)=\frac{\log s_j(z)}{2\pi i},\qquad
 h_j[z,\zeta]=[z,\zeta+\alpha(z)\Pi(z)v_j].
 \tag{FF36}\label{ffcert:eq:overlap-map}
\]
A branch change adds $n\Pi(z)v_j$, an original lattice period,
so the map is well-defined and holomorphic on the punctured torus
quotient; subtraction of the same term gives its inverse. Under the
clockwise step $s_j\mapsto\zeta_js_j$, the compatible branch is
$\alpha(g_jz)=\alpha(z)-1/m_j$. The covariance
\eqref{ffcert:eq:varying-action} and $A_jv_j=v_j$ therefore give
the exact conjugacy to the original linear lift
$\widehat G_j^0(z,\zeta)=(g_jz,R_{g_j}(z)\zeta)$:
\[
 h_j\widehat G_jh_j^{-1}=\widehat G_j^0.
 \tag{FF37}\label{ffcert:eq:overlap-conjugacy}
\]
Indeed the two translations in $h_j\widehat G_j$ add to
$\Pi(g_jz)v_j(1/m_j+\alpha(g_jz))
=\alpha(z)\Pi(g_jz)v_j$, exactly the translation in
$\widehat G_j^0h_j$.

The corresponding real translation retains the full complex structure:
\[
\begin{gathered}
 J_z=\Pi(z)^{-1}i\Pi(z),\qquad
 a_j(z)=\operatorname{Re}\alpha(z)v_j
               +\operatorname{Im}\alpha(z)J_zv_j,\\
 J_{g_jz}A_j=A_jJ_z,\qquad
 a_j(g_jz)=A_ja_j(z)-v_j/m_j.
\end{gathered}
\tag{FF38}\label{ffcert:eq:real-overlap}
\]
This follows by applying $\Pi(z)^{-1}$ to
$\alpha(z)\Pi(z)v_j$ and then using covariance; complex scalar
multiplication is not being imposed on the marked real lattice.
The inverse image of the middle zero section is consequently
$-a_j(z)$, whose endpoint after one clockwise step is exactly
$A_j(-a_j(z))+v_j/m_j$. Iterating $m_j$ such steps increases
this marked endpoint by precisely $v_j$, not by an unspecified
translated generator. Thus the overlap map realizes the based
peripheral convention in \eqref{ffcert:eq:whole-kernel}.

The map from the fibre over $z$ to its deck-translated fibre over
$g_jz$ has marked linear part $A_j$; its inverse has part
$A_j^{-1}$. The original contragredient convention is
\[
 (A_j^{-1})^{\mathsf t}=T_j,
 \qquad A_j=(T_j^{-1})^{\mathsf t}.
 \tag{FF39}\label{ffcert:eq:contragredient-direction}
\]
These directions keep a based loop distinct from its inverse deck
map. No matrix or generator is reversed to obtain the attachment relation.

\begin{corollary}[The exact integral first-cohomology covering image]
\label{ffcert:cor:covering-h1}
For the original covering $\pi_j:A_0^{(j)}\to S_j$,
\[
 H_1(S_j,\mathbb Z)\cong\mathbb Z^2,\qquad
 \pi_j^*H^1(S_j,\mathbb Z)
   =\mathbb Z(m_j\gamma)\oplus\mathbb Z\psi_j
        \subset H^1(A_0^{(j)},\mathbb Z).
 \tag{FF40}\label{ffcert:eq:integral-h1-image}
\]
The index $m_j$ in the invariant covector lattice is retained.
\end{corollary}

\begin{proof}
By \cref{ffcert:prop:integral-summands}, the map
\[
 \Lambda\longrightarrow\mathbb Z^2,\qquad
 \lambda\longmapsto(\gamma(\lambda),\psi_j(\lambda))
 \tag{FF41}\label{ffcert:eq:coinvariant-coordinates}
\]
has kernel exactly $(A_j-I)\Lambda$. It is onto:
$\widehat\gamma$ gives $(1,0)$, and $\widehat w$ for $j=1$,
or $\widehat u$ for $j=2$, gives $(0,1)$. Thus the full
coinvariant group is $\mathbb Z^2$. In these coordinates $v_j$
is $(1,0)$ or $(-1,0)$. Abelianizing
\eqref{ffcert:eq:full-deck-group} adjoins one generator with relation
$m_j\widehat G_j=v_j$, eliminating a primitive generator and
leaving $\mathbb Z^2$. This computes $H_1(S_j,\mathbb Z)$ without
global topology. A homomorphism from this group to $\mathbb Z$
restricts on $\Lambda$ to $a\gamma+b\psi_j$ precisely when
$m_j\mid a$, since it must assign the integer
$a\gamma(v_j)/m_j$ to $\widehat G_j$. This necessary and sufficient
condition proves the exact covering pullback.
\end{proof}

\subsubsection{The finite-to-global analytic gluing interface}

The global construction of \cref{thm:candidate-global-geometry} glues
complete punctured restrictions of the original middle family and
the three local fillings. The finite identifications are precisely
\eqref{ffcert:eq:overlap-map}. The special discs are pairwise disjoint,
and a middle point over each punctured disc has exactly one
representative in its filling. Thus an equivalence class meets any
given summand at most once. Saturations of open subsets are open,
because the overlap identifications are biholomorphisms, so each summand
embeds openly and injectively.

Two points at one special base value have representatives in its same
filling; two points at one ordinary base value have representatives in
the middle piece. They have disjoint neighbourhoods in that Hausdorff
open summand, and their images remain disjoint because no equivalence
class meets that summand twice. Points at distinct base values are
separated by the base. This proves the gluing's Hausdorff assertion
using the actual full restrictions, not an unrestricted gluing principle.

In particular $f^{-1}(D_j)$ is the entire corresponding filling:
every middle point over $D_j^*$ is identified with its unique filling
representative. Thus the properness in
\cref{ffcert:prop:quotient-charts} is properness of the full inverse
image. For a compact base set $K$, finitely many relatively compact
$V_a$ whose closures lie inside members $U_a$ of the original
four-piece base cover can be chosen to cover $K$. The compact sets
$K\cap\overline V_a$ have compact preimages under the corresponding
full proper restrictions, and these preimages cover exactly
$f^{-1}(K)$. This is the compact-cover argument for global properness.
The separate cusp and middle constructions supply their own full
proper restrictions; the finite-local proof neither replaces those
constructions nor implies smooth sphere recognition.

\subsubsection{The exact normal character and a smooth trivialization}

\begin{theorem}[Holomorphic normal order and the explicit smooth product]
\label{ffcert:thm:normal-trivialization}
Let $\mathcal N_j=N_{S_j/N_j}$, equivalently
$\mathcal O_X(S_j)|_{S_j}$ after the specified open embedding into
the original global analytic construction. This line has exact order
$m_j$ in the holomorphic Picard group, but has the following explicit
smooth trivialization and hence zero integral first Chern class.
\end{theorem}

\begin{proof}
Differentiating the actual transverse coordinate in
\eqref{ffcert:eq:disc-and-phase} identifies the normal line as the
quotient of the trivial line over $A_0^{(j)}$ by
\[
 (x,n)\longmapsto(A_jx+v_j/m_j,\zeta_j n),
 \qquad \zeta_j=e^{-2\pi i/m_j}.
 \tag{FF42}\label{ffcert:eq:normal-action}
\]
Here $x$ is the original real marking of the central complex torus;
the action on the torus and line is holomorphic. A section of its
$k$-th power therefore pulls back to a function satisfying
\[
 h(A_jx+v_j/m_j)=\zeta_j^k h(x).
 \tag{FF43}\label{ffcert:eq:normal-section-character}
\]
The positive exponent follows from the quotient itself: the representative
$(x,h(x))$ is identified at the image point with
$(A_jx+v_j/m_j,\zeta_j^kh(x))$, which must equal the representative
of the section there. The negative exponent would instead describe
the dual normal line under the declared action.

The $m_j$-th power has the trivial character, so the constant function
one descends to a nowhere-zero holomorphic section of that power.
If a power $0<k<m_j$ were holomorphically trivial, its trivializing
section would pull back to a nowhere-zero holomorphic function on the
compact connected torus $A_0^{(j)}$. Such a function is constant,
and \eqref{ffcert:eq:normal-section-character} would force
$\zeta_j^k=1$, impossible in this range. This proves the exact
holomorphic order in \cref{lem:finite-normal-order-audit}, with the
normal tangent character fully specified.

For the smooth trivialization put
\[
 \epsilon_j=\gamma(v_j)\in\{1,-1\},\qquad
 H_j(x)=\exp(-2\pi i\epsilon_j\gamma(x))
 \quad(x\in\Lambda\otimes\mathbb R).
 \tag{FF44}\label{ffcert:eq:smooth-normal-section}
\]
Integrality of $\gamma(\lambda)$ gives $H_j(x+\lambda)=H_j(x)$,
so this is a well-defined nowhere-zero smooth function on the actual
marked central torus. Moreover $\gamma A_j=\gamma$ and
$\epsilon_j^2=1$ give
\[
 \begin{aligned}
 H_j(A_jx+v_j/m_j)
 &=\exp(-2\pi i\epsilon_j\gamma(x))
       \exp(-2\pi i\epsilon_j^2/m_j)\\
 &=\zeta_jH_j(x).
 \end{aligned}
 \tag{FF45}\label{ffcert:eq:smooth-equivariance}
\]
It consequently descends to a nowhere-zero section of the actual
normal line. Explicitly the smooth fibrewise-linear trivialization and
its inverse are
\[
 \begin{aligned}
 \Theta_j:\mathcal N_j&\longrightarrow S_j\times\mathbb C,
 & [x,n]&\longmapsto([x],n/H_j(x)),\\
 \Theta_j^{-1}:S_j\times\mathbb C&\longrightarrow\mathcal N_j,
 & ([x],z)&\longmapsto[x,zH_j(x)].
 \end{aligned}
 \tag{FF46}\label{ffcert:eq:normal-product-isomorphism}
\]
Replacing $x$ by a lattice translate leaves both formulas unchanged.
Replacing $(x,n)$ by
$(A_jx+v_j/m_j,\zeta_jn)$ also leaves $n/H_j(x)$ unchanged by
\eqref{ffcert:eq:smooth-equivariance}; the inverse has the same
equivariance. The two maps are therefore well-defined, smooth, and
inverse, retaining the original signs for each $j$. Thus
$c_1(\mathcal N_j)=0$ in integral cohomology while its holomorphic
order is exactly $m_j$. This conclusion uses only the original affine
action and integral covector, not the topology of a global $X$.
\end{proof}
